\section{敏感性分析}

本章回答三个问题：哪些参数主导长期结果，多因素扰动下尾部风险有多大，外生风险变量和历史样本是否支持该判断。GPR、OVX、历史窗口和传统蒙特卡洛路径图放入附录，正文保留排序、区间和摘要证据。

\subsection{参数敏感性与关键驱动因素}

以中性情景为基准，对 16 个关键参数进行单因素扰动，观察第 180 天价格、外推期最高价和剩余供需缺口变化。设参数 \(\theta_j\) 在扰动集合 \(\Omega_j=\{\theta_j^{(1)},\ldots,\theta_j^{(m)}\}\) 上取值，其他参数固定。记长期递推得到的第 180 天价格、外推期最高价和第 180 天剩余缺口分别为 \(P_{180}(\theta_j^{(r)})\)、\(P_{\max}(\theta_j^{(r)})\) 和 \(G_{180}(\theta_j^{(r)})\)，则响应幅度定义为

\begin{equation}
\begin{aligned}
  \Delta P_{180}^{(j)}&=\max_{r}P_{180}(\theta_j^{(r)})-\min_{r}P_{180}(\theta_j^{(r)}),\\
  \Delta P_{\max}^{(j)}&=\max_{r}P_{\max}(\theta_j^{(r)})-\min_{r}P_{\max}(\theta_j^{(r)}),\\
  \Delta G_{180}^{(j)}&=\max_{r}G_{180}(\theta_j^{(r)})-\min_{r}G_{180}(\theta_j^{(r)}).
\end{aligned}
\end{equation}

综合敏感度定义为

\begin{equation}
  \mathcal S_j
  =
  \Delta P_{180}^{(j)}
  +0.45\Delta P_{\max}^{(j)}
  +0.002\Delta G_{180}^{(j)}.
\end{equation}

该指标兼顾终点价格、阶段峰值和供需风险；系数 \(0.45\) 和 \(0.002\) 将峰值风险和供需缺口映射到同一量级，使排序反映综合影响。

\begin{table}[!htbp]
\centering
\caption{参数综合敏感度排序}
\begin{tabularx}{\linewidth}{S[table-format=2.0] Y S[table-format=2.2] S[table-format=2.2] S[table-format=2.2]}
\toprule
{排名} & 参数 & {综合敏感度} & {第180天波动} & {外推峰值波动} \\
\midrule
1 & 封锁风险衰减速度 & 15.60 & 12.58 & 6.72 \\
2 & 地缘风险权重 & 13.36 & 8.26 & 11.33 \\
3 & 供应中断量 & 10.76 & 8.41 & 4.34 \\
4 & 不确定性与制度风险强度 & 9.25 & 3.95 & 11.77 \\
5 & 制度风险占比 & 7.89 & 5.40 & 5.53 \\
6 & 长期需求弹性 & 6.20 & 5.00 & 1.80 \\
7 & SPR 释放上限 & 5.04 & 4.31 & 1.08 \\
8 & 冲击不确定性占比 & 3.01 & 1.61 & 3.11 \\
9 & 信心恢复后风险底座 & 2.48 & 2.44 & 0.09 \\
10 & 绕道运输能力 & 0.66 & 0.00 & 1.47 \\
11 & 恐慌衰减速度 & 0.20 & 0.01 & 0.44 \\
12 & 预期修复强度 & 0.13 & 0.00 & 0.30 \\
13 & 过剩供给回归强度 & 0.10 & 0.08 & 0.05 \\
14 & SPR 启动延迟 & 0.07 & 0.00 & 0.15 \\
15 & 绕道恢复耗时 & 0.07 & 0.00 & 0.15 \\
16 & SPR 收缩起点 & 0.03 & 0.00 & 0.07 \\
\bottomrule
\end{tabularx}
\end{table}

\begin{figure}[H]
\centering
\includegraphics[width=0.86\linewidth]{参数敏感性龙卷风图.png}
\caption{关键参数综合敏感度排序}
\label{fig:sensitivity-tornado}
\end{figure}

\begin{figure}[H]
\centering
\includegraphics[width=0.86\linewidth]{长期敏感性扰动热力图.png}
\caption{长期敏感性扰动热力图}
\label{fig:sensitivity-heatmap}
\end{figure}

\FloatBarrier

敏感性结果表明，第 180 天价格对封锁风险衰减速度最敏感。当该速度从 0.002 扰动到 0.008 时，第 180 天价格波动达到 12.58 美元/桶。接入 SPR 有限预算后，供应中断量和 SPR 释放上限排序上升，说明长期结果同时受风险吸收、物理缺口和政策缓冲约束。

\keypoint{长期结果主要受封锁风险衰减速度、地缘风险权重、供应中断量和制度风险强度影响，需与 SPR 有限预算和尾部区间合读。}

\subsection{蒙特卡洛情景树与尾部风险}
\label{sec:monte-carlo}

单因素敏感性只给出边际影响，长期风险则常由多因素共同恶化触发。本文在赛题边界和扰动范围内进行 2000 次联合抽样。抽样使用压力指数 \(z\in[0,1]\)：\(z\) 越高，供应中断量、地缘风险权重和制度风险强度越高，SPR 可用释放量、绕道能力、需求调整速度和预期修复强度越低。每次抽样均调用同一套中长期油价调节模型。

扰动尺度由 2017--2025 年附件历史数据校准。共构造 2146 个同长度窗口；2026 冲突窗口的日收益波动率位于历史约 95.5\% 分位，最大单日跳变位于约 89.4\% 分位。据此设定压力指数噪声标准差为 0.073，状态转移日度扰动波动率为 0.0132。

\begin{table}[!htbp]
\centering
\caption{蒙特卡洛情景树尾部风险结果}
\begin{tabular}{lr}
\toprule
指标 & 数值 \\
\midrule
外推期最高价突破 120 美元/桶概率 & 17.3\% \\
外推期最高价突破 130 美元/桶概率 & 4.2\% \\
第 180 天剩余缺口超过 500 万桶/日概率 & 12.5\% \\
二次跳涨风险为高概率 & 13.3\% \\
第 180 天价格 P05 & 90.03 \\
第 180 天价格 P50 & 98.00 \\
第 180 天价格 P95 & 111.95 \\
外推期最高价 P95 & 128.84 \\
\bottomrule
\end{tabular}
\end{table}

\begin{figure}[H]
\centering
\includegraphics[width=0.88\linewidth,height=0.34\textheight,keepaspectratio]{蒙特卡洛尾部风险概率.png}
\caption{蒙特卡洛尾部风险概率}
\end{figure}

\FloatBarrier

蒙特卡洛路径第 180 天中位数为 98.00 美元/桶，与中性情景 94.47 美元/桶处在同一数量级。外推期最高价突破 120 美元/桶的概率为 17.3\%，突破 130 美元/桶的概率为 4.2\%。因此，长期风险重点不在中性点预测，而在多因素联动恶化下的尾部组合。

状态转移层令 \(Z_t\in\{E,N,U\}\) 在缓和、维持和升级之间切换，转移概率由剩余缺口、SPR 预算消耗、库存消耗、绕道瓶颈、高价压力、GPR 和 OVX 在 logit 空间修正。3000 条路径显示，第 180 天价格中位数为 97.54 美元/桶，5\%--95\% 区间为 91.27--112.76 美元/桶；外推期最高价 P95 为 134.55 美元/桶，突破 120 美元/桶的条件概率为 76.3\%。

\subsection{外生风险变量与历史稳健性摘要}

GPR 和 OVX 不替代附件真实价格，也不作为短期同日拟合变量。GPR 刻画地缘风险叙事，OVX 反映期权市场对未来波动的定价；二者用于约束长期状态压力和尾部风险。

\begin{table}[H]
\centering
\caption{外生风险变量与历史稳健性摘要}
\begin{tabularx}{\linewidth}{p{0.20\linewidth}p{0.28\linewidth}X}
\toprule
证据来源 & 关键结果 & 模型含义 \\
\midrule
GPR 月度指数 & 2026 年 3 月历史分位数 98.98\%，4 月历史分位数 97.76\%；滞后 1 月 GPR 与本月收益相关系数仅 0.030 & 地缘风险处于罕见高位，但不应包装成强短期收益预测器；适合进入长期风险压力约束。 \\
OVX 隐含波动率 & 冲突窗口均值 87.07，最大值 120.91，均值历史分位 97.3\%；滞后 1 日 OVX 与 7 日实现波动率相关系数 0.790 & 期权市场确认冲突窗口的高波动风险；适合约束长期波动和尾部概率，而不直接预测方向。 \\
历史同长度窗口 & 2026 冲突窗口累计收益率历史分位 99.3\%，实现波动率 95.7\%，朴素基准 RMSE 98.5\% & 说明该窗口确实是高涨幅、高波动、高预测难度样本；历史数据用于背景分布和扰动尺度校准。 \\
\bottomrule
\end{tabularx}
\end{table}

\FloatBarrier

上述结果给出三点判断：外生变量支持“冲突窗口风险读数异常高”的前提；GPR 和 OVX 对方向收益的直接解释力有限，不宜作为短期点预测变量；完整历史数据用于校准波动尺度、基准难度和极端性位置，不替代霍尔木兹封锁机制。
